\section{Part II - Multivariable control}
\begin{itemize}
    \item Introduce and motivate the task. You can include a brief summary of the different tasks in the lab-preparation assignment which was necessary to complete the lab. No extensive listing or derivation of equations are needed (you can mention, e.g., LQR, controllability, integral effect, augmented matrices etc.). 
    \item Include your hypothesis/test-plan for the matrices $\mathbf{Q}_{\text{LQR}}$ and $\mathbf{R}_{\text{LQR}}$ both with and without integral effect and explain your thought process behind it. Include it if you did any experiments with the values in $\mathbf{F}$.
    \item Include the results from the experiments that were done in the lab according to the test-plan, and any other tests you found interesting during your time in the lab. Include plots and discuss the result. This is the most important part of the report. How you structure this part of the report is up to you. 
    \item Conclusion. Include a short paragraph summarizing the results.
\end{itemize}

The purpose of this task is to control the helicopter using LQR control. To do so, we first need to find the state space representation.

\begin{align*}
    \dot{x}&=Ax+Bu \\
    x&=
    \begin{bmatrix}
        p \\
        \dot{p} \\
        \dot{e} \\
    \end{bmatrix}
    \text{ and } u=
    \begin{bmatrix}
        \tilde{V_s} \\
        V_d
    \end{bmatrix}\\
    \ddot{p}&=K_1 V_d \\
    \ddot{e}&=K_2 \tilde{V_s}\\
    \dot{x}&=
    \begin{bmatrix}
        0&1&0 \\
        0&0&0 \\
        0&0&0 \\
    \end{bmatrix}
    x+
    \begin{bmatrix}
        0&0 \\
        0&K_1 \\
        K_2&0
    \end{bmatrix}
    \begin{bmatrix}
        \tilde{V_s} \\
        V_d
    \end{bmatrix}\\
\end{align*}


Then we need to make sure that 
\begin{equation*}
    \lim_{t \to \infty} \implies \dot{x} = 0 \text{ and } x = \left[
    \begin{array}{c}
        p_c \\
        0 \\
        \dot{e}_c
    \end{array}
    \right]
\end{equation*}
This gives the equation
\begin{align*}
    \dot{x} &= Ax + Bu \\
    &= (A - BK)x + BFr \\
    &= 0
\end{align*}
And solved for F we get
\begin{equation*}
    F =
    \begin{bmatrix}
        k_{11} & k_{13} \\
        k_{21} & k_{23}
    \end{bmatrix}
\end{equation*}

Such that
\begin{equation*}
    \lim_{t \to \infty} 
    \begin{bmatrix}
        p \\
        \dot{e}
    \end{bmatrix}
    =
    \begin{bmatrix}
        p_c \\
        \dot{e_c}
    \end{bmatrix}
\end{equation*}

\subsection{Hypotheses and test plan}
A high q/r will give a fast response, but with a high energy drainage. Can expect some overswing as that is a waste of energy. A low q/r will give a slow response where you minimize the amount of energy spent. We also think that high isolated Q values regulates that parameter better, at the cost of the others. We also think that implemented integral action will remove standard deviaton from the system.

This is our overarching test plan:
\begin{itemize}
    \item Try 10 on all diagonals
    \item try high q
    \item try high r
    \item try different values inside of q and r
\end{itemize}
The specific test plan for the 3-state system is shown in \autoref{tab:3stateHyp}, and the 5-state is shown in \autoref{tab:5stateHyp}.
\begin{table}[htbp]
    \centering
    \caption{Test values for 3 state LQR}
    \begin{tabular}{c|ccccc}
        \textbf{Test} & \( Q_p \) & \( Q_{\dot{p}} \) & \( Q_{\dot{e}} \) & \( R_{\tilde{V}_s} \) & \( R_{V_d} \) \\
        \hline
        1  & 10  & 10  & 10  & 10  & 10  \\
        2  & 200 & 10  & 10  & 10  & 10  \\
        3  & 10  & 200 & 10  & 10  & 10  \\
        4  & 10  & 10  & 200 & 10  & 10  \\
        5  & 10  & 10  & 10  & 200 & 10  \\
        6  & 10  & 10  & 10  & 10  & 200 \\
        7  & 200 & 200 & 200 & 10  & 10  \\
        8 & 10  & 10  & 10  & 200 & 200 \\
    \end{tabular}
    \label{tab:3stateHyp}
\end{table}

\begin{table}[htbp]
    \centering
    \caption{Test values for 5 state LQR}
    \begin{tabular}{c|ccccccc}
        \textbf{Test} & \( Q_p \) & \( Q_{\dot{p}} \) & \( Q_{\dot{e}} \) & \( Q_{\dot{\gamma}} \) & \( Q_{\dot{\zeta}} \) & \( R_{\tilde{V}_s} \) & \( R_{V_d} \) \\
        \hline
        1  & 10  & 10  & 10  & 10  & 10  & 10  & 10  \\
        2  & 200 & 10  & 10  & 10  & 10  & 10  & 10  \\
        3  & 10  & 200 & 10  & 10  & 10  & 10  & 10  \\
        4  & 10  & 10  & 200 & 10  & 10  & 10  & 10  \\
        5  & 10  & 10  & 10  & 200 & 10  & 10  & 10  \\
        6  & 10  & 10  & 10  & 10  & 200 & 10  & 10  \\
        7  & 10  & 10  & 10  & 10  & 10  & 200 & 10  \\
        8  & 10  & 10  & 10  & 10  & 10  & 10  & 200 \\
        9  & 200 & 200 & 200 & 200 & 200 & 10  & 10  \\
        10 & 10  & 10  & 10  & 10  & 10  & 200 & 200 \\
    \end{tabular}
    \label{tab:5stateHyp}
\end{table}
% Helper function to plot subfigures fast
% #1 = path prefix, #2 = number of figures
\newcommand{\insertsubfigs}[2]{%
  \foreach \i in {1,...,#2}{%
    \begin{subfigure}{0.40\textwidth}
      \includegraphics[width=\textwidth]{#1\i.eps}
      \caption{Results for test \i.}
    \end{subfigure}
  }%
}

\subsection{Results}
From \autoref{fig:LQR_3state}, we see that there is a lot of standard deviation in the 3 state system. We see overall better results, without standard deviation, for the 5-state regulator (\autoref{fig:LQR_5state}), as expected. In the 5-state system, the results from test 2 in \autoref{fig:LQR_5state}. were the best, which makes sense since we tested pitch, and the cost of bad pitch was proportionally higher than everything else in test 2. We also see that test 9 was better than a lot of the others with high Q, which confirmes that to increase the cost of one state, affects the other states. For instance test 3 was quite bad for the pitch.

\begin{figure}[h!]
    \centering
    \insertsubfigs{figures/labday2/3state/plot_}{8}
    \caption{Step response for the three-state LQR regulator.}
    \label{fig:LQR_3state}
\end{figure}



\begin{figure}[h]
    \centering
    \insertsubfigs{figures/labday2/5state/plot_}{10}
    \caption{Step response for the five-state LQR regulator.}
    \label{fig:LQR_5state}
\end{figure}